\begin{proof}[Proof of \cref{obj:thm:three-output-region}]
\begin{enumerate}
\item \textit{Certificate notation.}
Write the parameter domain and the active pattern set as
\[
\mathcal D:=(0,1)^3\times(0,\infty),
\qquad
\mathcal K:=\{S_1,S_6,S_8\}.
\]
For a parameter vector, use the coordinate notation
\[
\xi=(p(\xi),\mu_0(\xi),\mu_1(\xi),\varepsilon(\xi)),
\qquad
\theta(\xi):=(\mu_0(\xi),\mu_1(\xi))^\top.
\]
The fixed pattern enumeration is
\[
\mathcal S:=\{S_k:1\le k\le14\},
\qquad
S_k:=\{j\in\{0,1,2,3\}:\text{the \(j\)-th binary digit of \(k\) is \(1\)}\}.
\]
In particular,
\[
S_1=\{0\},\qquad S_6=\{1,2\},\qquad S_8=\{3\}.
\]
For \(S\in\mathcal S\), define its staircase ray by
\[
b_{S,j}(\varepsilon):=
\begin{cases}
e^\varepsilon,&j\in S,\\
1,&j\notin S,
\end{cases}
\qquad j\in\{0,1,2,3\},
\qquad
b_S(\varepsilon):=
(b_{S,0}(\varepsilon),\ldots,b_{S,3}(\varepsilon))^\top.
\]
The record probabilities, pattern gradients, and pattern masses are
\[
\pi_\theta(p):=
\begin{pmatrix}
(1-p)(1-\mu_0)\\
(1-p)\mu_0\\
p(1-\mu_1)\\
p\mu_1
\end{pmatrix},
\qquad
g_S(p,\varepsilon):=
\begin{pmatrix}
(1-p)(b_{S,1}(\varepsilon)-b_{S,0}(\varepsilon))\\
p(b_{S,3}(\varepsilon)-b_{S,2}(\varepsilon))
\end{pmatrix},
\]
\[
h_S(\theta;p,\varepsilon):=
\sum_{j=0}^3\pi_{\theta,j}(p)b_{S,j}(\varepsilon),
\qquad
v(t):=(t,t+1)^\top,\quad t\in\mathbb R.
\]
The additional arguments on \(h_S\) display its dependence on assignment and privacy. For \(\xi\in\mathcal D\), define the pattern information and its profiling slope by
\[
f_S(\xi,t):=
\frac{\bigl(g_S(p(\xi),\varepsilon(\xi))^\top v(t)\bigr)^2}
{h_S(\theta(\xi);p(\xi),\varepsilon(\xi))},
\]
\[
\sigma_S(\xi,t):=
\frac{
2\bigl(g_S(p(\xi),\varepsilon(\xi))^\top v(t)\bigr)
\bigl(g_{S,0}(p(\xi),\varepsilon(\xi))
      +g_{S,1}(p(\xi),\varepsilon(\xi))\bigr)}
{h_S(\theta(\xi);p(\xi),\varepsilon(\xi))}.
\]
The masses are positive on this domain, as shown below, so these quotients are well defined. The feasible set and objective are
\[
\mathcal A_\varepsilon:=
\left\{\beta\in[0,\infty)^{14}:
\sum_{S\in\mathcal S}\beta_S b_{S,j}(\varepsilon)=1
\text{ for every }j\in\{0,1,2,3\}\right\},
\]
\[
F_{\theta(\xi)}(\beta,t)
=\sum_{S\in\mathcal S}\beta_S f_S(\xi,t),
\]
where the objective in this display is evaluated at \(p(\xi)\) and \(\varepsilon(\xi)\). These are the quantities entering \cref{obj:def:staircase-saddle}. The three-pattern certificate of \cref{obj:def:support-regions} has support \(\mathcal K\), positive active weights, zero weighted profiling slope, active dual equalities, and strict inactive dual inequalities.
% lean: CausalSmith.Stat.LdpAteEfficiencySurface.certifiedRegion

\item \textit{Canonical weights for an arbitrary certificate.}
Fix \(\xi_0\in\mathcal R_3\), and choose witnessing quantities
\[
\beta^{(0)}\in\mathcal A_{\varepsilon(\xi_0)},
\qquad t_{\mathrm{wit}}\in\mathbb R,
\qquad \eta^{(0)}\in\mathbb R^4.
\]
Their support is \(\mathcal K\), so
\[
\beta^{(0)}_S=0\quad(S\notin\mathcal K).
\]
Put
\[
r_0:=e^{\varepsilon(\xi_0)}>1.
\]
Normalization at input coordinates \(0,1,3\) gives
\[
\begin{aligned}
r_0\beta^{(0)}_{S_1}+\beta^{(0)}_{S_6}
 +\beta^{(0)}_{S_8}&=1,\\
\beta^{(0)}_{S_1}+r_0\beta^{(0)}_{S_6}
 +\beta^{(0)}_{S_8}&=1,\\
\beta^{(0)}_{S_1}+\beta^{(0)}_{S_6}
 +r_0\beta^{(0)}_{S_8}&=1.
\end{aligned}
\]
Subtracting consecutive equations yields
\[
(r_0-1)(\beta^{(0)}_{S_1}-\beta^{(0)}_{S_6})=0,
\qquad
(r_0-1)(\beta^{(0)}_{S_6}-\beta^{(0)}_{S_8})=0.
\]
Consequently,
\[
\beta^{(0)}=\alpha^{(3)}(\varepsilon(\xi_0)),
\qquad
\alpha^{(3)}_S(\varepsilon):=
\begin{cases}
(e^\varepsilon+2)^{-1},&S\in\mathcal K,\\
0,&S\notin\mathcal K,
\end{cases}
\quad \varepsilon>0.
\]
For every \(\varepsilon>0\), the active weights are positive. Since the three active patterns partition the inputs,
\[
\sum_{S\in\mathcal K}b_{S,j}(\varepsilon)=e^\varepsilon+2
\quad(j=0,1,2,3).
\]
Thus
\[
\alpha^{(3)}(\varepsilon)\in\mathcal A_\varepsilon,
\qquad
\{S:\alpha^{(3)}_S(\varepsilon)\ne0\}=\mathcal K.
\]
% lean: CausalSmith.Stat.LdpAteEfficiencySurface.threeMaskWeights_unique
% lean: CausalSmith.Stat.LdpAteEfficiencySurface.threeMaskWeights_feasible
% lean: CausalSmith.Stat.LdpAteEfficiencySurface.threeMaskWeights_support

\item \textit{A continuous stationary direction.}
For every \(\xi\in\mathcal D\), all record probabilities are nonnegative and all staircase-ray entries are positive. In particular,
\[
h_S(\theta(\xi);p(\xi),\varepsilon(\xi))
\ge
(1-p(\xi))(1-\mu_0(\xi))b_{S,0}(\varepsilon(\xi))
>0
\quad(S\in\mathcal S).
\]
Define the canonical stationarity function and its linear coefficient by
\[
L_3(\xi,t):=
\sum_{S\in\mathcal S}
\alpha^{(3)}_S(\varepsilon(\xi))\sigma_S(\xi,t),
\qquad
a_3(\xi):=L_3(\xi,1)-L_3(\xi,0).
\]
The identity
\[
g_S^\top v(t)=g_{S,1}+t(g_{S,0}+g_{S,1})
\]
gives, with all quantities evaluated at \(\xi\),
\[
\sigma_S(\xi,t)
=
\sigma_S(\xi,0)
+
t\,\frac{2(g_{S,0}+g_{S,1})^2}{h_S(\theta(\xi))}.
\]
Hence
\[
L_3(\xi,t)=L_3(\xi,0)+t\,a_3(\xi),
\]
\[
a_3(\xi)=
\sum_{S\in\mathcal S}
\alpha^{(3)}_S(\varepsilon(\xi))
\frac{2(g_{S,0}+g_{S,1})^2}{h_S(\theta(\xi))}>0.
\]
Every summand is nonnegative, and the \(S_1\) summand is strictly positive because
\[
g_{S_1,0}+g_{S_1,1}
=-(e^{\varepsilon(\xi)}-1)(1-p(\xi))\ne0.
\]
We may therefore define, throughout \(\mathcal D\),
\[
t_3(\xi):=-\frac{L_3(\xi,0)}{a_3(\xi)},
\qquad
L_3(\xi,t_3(\xi))=0.
\]
The witnessing stationarity equation and the canonical-weight identity give
\[
0=L_3(\xi_0,t_{\mathrm{wit}})
=L_3(\xi_0,0)+t_{\mathrm{wit}}a_3(\xi_0),
\qquad
t_{\mathrm{wit}}=t_3(\xi_0).
\]

For each fixed pattern and input coordinate, the ray entry is either \(e^\varepsilon\) or \(1\), and is continuous. The masses are finite sums of products of continuous functions; the gradients are continuous as well. Their positive denominators show that \(L_3(\xi,0)\) and \(L_3(\xi,1)\) are continuous at \(\xi_0\). Since \(a_3(\xi_0)>0\), the quotient defining \(t_3\) is continuous there. It follows that
\[
\xi\longmapsto f_S(\xi,t_3(\xi))
\quad\text{is continuous at }\xi_0
\quad(S\in\mathcal S).
\]
% lean: CausalSmith.Stat.LdpAteEfficiencySurface.patternMass_pos_of_interior
% lean: CausalSmith.Stat.LdpAteEfficiencySurface.r3StationarityCoeff_pos
% lean: CausalSmith.Stat.LdpAteEfficiencySurface.r3Stationarity_affine
% lean: CausalSmith.Stat.LdpAteEfficiencySurface.continuousAt_r3Direction_region

\item \textit{A local dual solution anchored at the witness.}
Define the active dual matrix and target vector on \(\mathcal D\) by
\[
r(\xi):=e^{\varepsilon(\xi)},
\qquad
\mathsf A_3(\xi):=
\begin{pmatrix}
r(\xi)&1&1&1\\
1&r(\xi)&r(\xi)&1\\
1&1&1&r(\xi)
\end{pmatrix},
\qquad
\boldsymbol f_3(\xi):=
\begin{pmatrix}
f_{S_1}(\xi,t_3(\xi))\\
f_{S_6}(\xi,t_3(\xi))\\
f_{S_8}(\xi,t_3(\xi))
\end{pmatrix}.
\]
The active equalities at \(\xi_0\), together with
\(t_{\mathrm{wit}}=t_3(\xi_0)\), say
\[
\mathsf A_3(\xi_0)\eta^{(0)}=\boldsymbol f_3(\xi_0).
\]
Use the fixed right inverse
\[
\mathsf R_0:=
\frac{1}{(r_0-1)(r_0+2)}
\begin{pmatrix}
r_0+1&-1&-1\\
-1&r_0+1&-1\\
0&0&0\\
-1&-1&r_0+1
\end{pmatrix}.
\]
Its denominator is nonzero. Multiplication gives
\[
\mathsf A_3(\xi_0)\mathsf R_0
=
\frac{1}{(r_0-1)(r_0+2)}
\begin{pmatrix}
r_0(r_0+1)-2&0&0\\
0&r_0(r_0+1)-2&0\\
0&0&r_0(r_0+1)-2
\end{pmatrix}
=\operatorname{diag}(1,1,1).
\]
The matrix \(\mathsf A_3\) and the target \(\boldsymbol f_3\) are continuous at \(\xi_0\). Since
\[
\det(\mathsf A_3(\xi_0)\mathsf R_0)=1,
\]
continuity of the determinant gives a neighborhood of \(\xi_0\) within
\[
\mathcal U_{\mathrm{inv}}
:=
\{\xi\in\mathcal D:
\det(\mathsf A_3(\xi)\mathsf R_0)\ne0\}.
\]
On this set define
\[
\eta_{\mathrm{loc}}(\xi):=
\eta^{(0)}
+
\mathsf R_0
(\mathsf A_3(\xi)\mathsf R_0)^{-1}
\bigl(\boldsymbol f_3(\xi)-\mathsf A_3(\xi)\eta^{(0)}\bigr).
\]
The inverse in this formula is taken on the stated invertible domain. Continuity of matrix inversion and the displayed formula imply continuity of \(\eta_{\mathrm{loc}}\) at \(\xi_0\). The residual vanishes at the base point, so
\[
\eta_{\mathrm{loc}}(\xi_0)=\eta^{(0)}.
\]
Moreover, for every \(\xi\in\mathcal U_{\mathrm{inv}}\),
\[
\begin{aligned}
\mathsf A_3(\xi)\eta_{\mathrm{loc}}(\xi)
&=
\mathsf A_3(\xi)\eta^{(0)}
+
\mathsf A_3(\xi)\mathsf R_0
(\mathsf A_3(\xi)\mathsf R_0)^{-1}
\bigl(\boldsymbol f_3(\xi)-\mathsf A_3(\xi)\eta^{(0)}\bigr)\\
&=\boldsymbol f_3(\xi).
\end{aligned}
\]
Thus the three active dual equalities hold locally, and the original inactive strict inequalities hold at the base point:
\[
f_S(\xi_0,t_3(\xi_0))
<
b_S(\varepsilon(\xi_0))^\top\eta_{\mathrm{loc}}(\xi_0)
\quad(S\notin\mathcal K).
\]
% lean: CausalSmith.Stat.LdpAteEfficiencySurface.r3DualMatrix_mul_rightInverse
% lean: CausalSmith.Stat.LdpAteEfficiencySurface.exists_eventually_anchoredSolution_of_rightInverse
% lean: CausalSmith.Stat.LdpAteEfficiencySurface.exists_local_r3Dual_of_mem

\item \textit{Persistence of the certificate.}
For every \(S\in\mathcal S\), the dual value
\[
b_S(\varepsilon(\xi))^\top\eta_{\mathrm{loc}}(\xi)
=
\sum_{j=0}^3b_{S,j}(\varepsilon(\xi))
\eta_{\mathrm{loc},j}(\xi)
\]
is continuous at \(\xi_0\), being a finite sum of products of continuous functions. Consequently, each inactive strict inequality persists in a neighborhood of \(\xi_0\). Intersect these finitely many neighborhoods with the neighborhood on which the active equations hold. The strict coordinate inequalities
\[
0<p(\xi)<1,\qquad
0<\mu_0(\xi)<1,\qquad
0<\mu_1(\xi)<1,\qquad
0<\varepsilon(\xi)
\]
also persist locally. We obtain a neighborhood \(\mathcal N\) of \(\xi_0\) such that, for every \(\xi\in\mathcal N\),
\[
\begin{gathered}
\xi\in\mathcal D,\qquad
\alpha^{(3)}(\varepsilon(\xi))\in\mathcal A_{\varepsilon(\xi)},\\
\{S:\alpha^{(3)}_S(\varepsilon(\xi))\ne0\}=\mathcal K,
\qquad
\alpha^{(3)}_S(\varepsilon(\xi))>0\quad(S\in\mathcal K),\\
\sum_{S\in\mathcal S}
\alpha^{(3)}_S(\varepsilon(\xi))\sigma_S(\xi,t_3(\xi))=0,\\
b_S(\varepsilon(\xi))^\top\eta_{\mathrm{loc}}(\xi)
=f_S(\xi,t_3(\xi))
\quad(S\in\mathcal K),\\
f_S(\xi,t_3(\xi))
<
b_S(\varepsilon(\xi))^\top\eta_{\mathrm{loc}}(\xi)
\quad(S\notin\mathcal K).
\end{gathered}
\]
These quantities constitute the certificate in \cref{obj:def:support-regions}, so
\[
\mathcal N\subseteq\mathcal R_3.
\]
As \(\xi_0\) was arbitrary, \(\mathcal R_3\) is open in \(\mathbb R^4\).
% lean: CausalSmith.Stat.LdpAteEfficiencySurface.isOpen_R3_continuousCertificate

\item \textit{Primal feasibility and stationarity at the stated point.}
Define
\[
\xi_{\mathrm{low}}
:=
\left(\frac3{10},\frac1{10},\frac1{10},\log3\right),
\qquad
\theta_{\mathrm{low}}:=\theta(\xi_{\mathrm{low}})
=\left(\frac1{10},\frac1{10}\right)^\top,
\qquad
t_{\mathrm{low}}:=-\frac{339}{9985},
\]
\[
\alpha:=\alpha^{(3)}(\log3),
\qquad
\alpha_S=
\begin{cases}
1/5,&S\in\mathcal K,\\
0,&S\notin\mathcal K.
\end{cases}
\]
Since \(\log3>0\) and \(e^{\log3}=3\), the canonical-weight calculation yields
\[
\alpha\in\mathcal A_{\log3},
\qquad
\{S:\alpha_S\ne0\}=\mathcal K,
\qquad
\alpha_{S_1}=\alpha_{S_6}=\alpha_{S_8}=\frac15.
\]
The record probabilities are
\[
\pi_{\theta_{\mathrm{low}}}(3/10)
=\frac1{100}(63,7,27,3)^\top.
\]
Substitution into the pattern formulas gives
\[
\begin{array}{c|c|c}
S&h_S(\theta_{\mathrm{low}})&g_S\\ \hline
S_1&113/50&(-7/5,0)^\top\\
S_6&42/25&(7/5,-3/5)^\top\\
S_8&53/50&(0,3/5)^\top
\end{array}
\]
and hence, for every \(t\in\mathbb R\),
\[
F_{\theta_{\mathrm{low}}}(\alpha,t)
=
\frac15\left[
\frac{98}{113}t^2
+\frac{(4t-3)^2}{42}
+\frac{18}{53}(t+1)^2
\right].
\]
Here and below objectives at \(\theta_{\mathrm{low}}\) are evaluated at \(p=3/10\) and \(\varepsilon=\log3\). The weighted slope at \(t_{\mathrm{low}}\) is therefore
\[
\sum_{S\in\mathcal S}\alpha_S
\sigma_S(\xi_{\mathrm{low}},t_{\mathrm{low}})
=
\frac15\left[
\frac{196}{113}t_{\mathrm{low}}
+\frac{8(4t_{\mathrm{low}}-3)}{42}
+\frac{36}{53}(t_{\mathrm{low}}+1)
\right]
=0.
\]
This establishes the required profiling stationarity.
% lean: CausalSmith.Stat.LdpAteEfficiencySurface.lowMeanCandidate_feasible
% lean: CausalSmith.Stat.LdpAteEfficiencySurface.lowMeanCandidate_stationary

\item \textit{The exact dual certificate and region membership.}
Define the four-coordinate dual vector by
\[
\eta_{\mathrm{low}}:=
\begin{pmatrix}
-21817593/398800900\\
32820291/8773619800\\
509870781/8773619800\\
41183667/398800900
\end{pmatrix}.
\]
Its coordinate sum is
\[
\sum_{j=0}^3\eta_{\mathrm{low},j}
=
-\frac{21817593}{398800900}
+\frac{32820291}{8773619800}
+\frac{509870781}{8773619800}
+\frac{41183667}{398800900}
=\frac{441}{3994}.
\]
For every \(S\in\mathcal S\), define its slack and the specialized ray entries by
\[
\Delta_S:=
b_S(\log3)^\top\eta_{\mathrm{low}}
-f_S(\xi_{\mathrm{low}},t_{\mathrm{low}}),
\qquad
b^{\mathrm{low}}_{S,j}:=b_{S,j}(\log3)
=
\begin{cases}
3,&j\in S,\\
1,&j\notin S.
\end{cases}
\]
The slack is computed directly as
\[
\Delta_S
=
\sum_{j=0}^3 b^{\mathrm{low}}_{S,j}\eta_{\mathrm{low},j}
-
\frac{
\left[
\frac7{10}(b^{\mathrm{low}}_{S,1}-b^{\mathrm{low}}_{S,0})t_{\mathrm{low}}
+
\frac3{10}(b^{\mathrm{low}}_{S,3}-b^{\mathrm{low}}_{S,2})(t_{\mathrm{low}}+1)
\right]^2}
{
\bigl(
63b^{\mathrm{low}}_{S,0}
+7b^{\mathrm{low}}_{S,1}
+27b^{\mathrm{low}}_{S,2}
+3b^{\mathrm{low}}_{S,3}
\bigr)/100
}.
\]
For example, \(S_2=\{1\}\) has ray \((1,3,1,1)^\top\), so
\[
\Delta_{S_2}
=
\frac{441}{3994}
+2\frac{32820291}{8773619800}
-\frac{98}{57}\left(\frac{339}{9985}\right)^2
=
\frac{1932296079}{16669877620}.
\]
Applying the same displayed formula to each of the fourteen binary patterns gives
\[
\begin{array}{c|c@{\qquad}c|c}
S&\Delta_S&S&\Delta_S\\ \hline
S_1&0&S_8&0\\
S_2&1932296079/16669877620&S_9&4394565/115652261\\
S_3&7441119/877361980&S_{10}&77641557/877361980\\
S_4&7441119/877361980&S_{11}&2821826721/35971841180\\
S_5&14134743/877361980&S_{12}&380056761/877361980\\
S_6&0&S_{13}&3683761767/11405705740\\
S_7&41610/3988009&S_{14}&50813490/115652261
\end{array}
\]
Thus the active and inactive cases satisfy, respectively,
\[
b_S(\log3)^\top\eta_{\mathrm{low}}
=f_S(\xi_{\mathrm{low}},t_{\mathrm{low}})
\quad(S\in\mathcal K),
\]
\[
f_S(\xi_{\mathrm{low}},t_{\mathrm{low}})
<
b_S(\log3)^\top\eta_{\mathrm{low}}
\quad(S\notin\mathcal K).
\]
The point \(\xi_{\mathrm{low}}\) belongs to \(\mathcal D\). Together with the feasible positive active weights and stationarity established above, these equalities and strict inequalities supply its certificate in \cref{obj:def:support-regions}. Therefore
\[
\xi_{\mathrm{low}}\in\mathcal R_3.
\]
% lean: CausalSmith.Stat.LdpAteEfficiencySurface.lowMeanDual_value
% lean: CausalSmith.Stat.LdpAteEfficiencySurface.lowMeanDual_slack
% lean: CausalSmith.Stat.LdpAteEfficiencySurface.lowMeanPoint_mem_R3

\item \textit{A uniform upper bound on profiled information.}
Let \(\beta\in\mathcal A_{\log3}\). Positive pattern masses and nonnegative weights imply
\[
F_{\theta_{\mathrm{low}}}(\beta,t)
=
\sum_{S\in\mathcal S}\beta_S
\frac{(g_S^\top v(t))^2}{h_S(\theta_{\mathrm{low}})}
\ge0
\quad(t\in\mathbb R).
\]
Thus its set of objective values,
\[
\mathcal U_\beta:=
\{F_{\theta_{\mathrm{low}}}(\beta,t):t\in\mathbb R\},
\]
is nonempty and bounded below by zero. In particular,
\[
\inf_{t\in\mathbb R}F_{\theta_{\mathrm{low}}}(\beta,t)
\le F_{\theta_{\mathrm{low}}}(\beta,t_{\mathrm{low}}).
\]
The active equality and inactive strict-inequality cases both give
\[
f_S(\xi_{\mathrm{low}},t_{\mathrm{low}})
\le b_S(\log3)^\top\eta_{\mathrm{low}}
\quad(S\in\mathcal S).
\]
Multiplication by \(\beta_S\ge0\) and summation yield
\[
F_{\theta_{\mathrm{low}}}(\beta,t_{\mathrm{low}})
\le
\sum_{S\in\mathcal S}\beta_S b_S(\log3)^\top\eta_{\mathrm{low}}.
\]
Normalization identifies the latter sum:
\[
\begin{aligned}
\sum_{S\in\mathcal S}\beta_S b_S(\log3)^\top\eta_{\mathrm{low}}
&=
\sum_{j=0}^3\eta_{\mathrm{low},j}
\left(\sum_{S\in\mathcal S}\beta_S b_{S,j}(\log3)\right)\\
&=\sum_{j=0}^3\eta_{\mathrm{low},j}
=\frac{441}{3994}.
\end{aligned}
\]
Consequently, for every feasible \(\beta\),
\[
\inf_{t\in\mathbb R}F_{\theta_{\mathrm{low}}}(\beta,t)
\le
F_{\theta_{\mathrm{low}}}(\beta,t_{\mathrm{low}})
\le\frac{441}{3994}.
\]
% lean: CausalSmith.Stat.LdpAteEfficiencySurface.lowMean_information_nonneg
% lean: CausalSmith.Stat.LdpAteEfficiencySurface.feasible_dualRay_sum
% lean: CausalSmith.Stat.LdpAteEfficiencySurface.lowMean_dual_upper_bound

\item \textit{The candidate infimum and optimized value.}
Expanding the candidate objective gives
\[
F_{\theta_{\mathrm{low}}}(\alpha,t)
=
\frac{39940}{125769}t^2+\frac8{371}t+\frac{411}{3710}
=
\frac{441}{3994}
+
\frac{39940}{125769}
\left(t+\frac{339}{9985}\right)^2.
\]
The coefficient of the square is positive. Therefore
\[
F_{\theta_{\mathrm{low}}}(\alpha,t)\ge\frac{441}{3994}
\quad(t\in\mathbb R),
\qquad
F_{\theta_{\mathrm{low}}}(\alpha,t_{\mathrm{low}})
=\frac{441}{3994},
\]
and hence
\[
\inf_{t\in\mathbb R}F_{\theta_{\mathrm{low}}}(\alpha,t)
=\frac{441}{3994}.
\]
Define the set of feasible profiled values by
\[
\mathcal E_{\mathrm{prof}}:=
\left\{
\inf_{t\in\mathbb R}F_{\theta_{\mathrm{low}}}(\beta,t):
\beta\in\mathcal A_{\log3}
\right\}.
\]
The uniform bound gives
\[
z\le\frac{441}{3994}\quad(z\in\mathcal E_{\mathrm{prof}}),
\]
while the feasible candidate gives
\[
\frac{441}{3994}\in\mathcal E_{\mathrm{prof}}.
\]
Thus this nonempty, bounded-above set has supremum \(441/3994\). By the optimized-value definition in \cref{obj:def:staircase-saddle},
\[
J^*(\theta_{\mathrm{low}},3/10,\log3)
=
\sup\mathcal E_{\mathrm{prof}}
=
\frac{441}{3994}
=
\inf_{t\in\mathbb R}F_{\theta_{\mathrm{low}}}(\alpha,t).
\]
This proves attainment by the stated weight vector.
% lean: CausalSmith.Stat.LdpAteEfficiencySurface.lowMeanCandidate_information_value
% lean: CausalSmith.Stat.LdpAteEfficiencySurface.lowMeanCandidate_global_min
% lean: CausalSmith.Stat.LdpAteEfficiencySurface.lowMeanCandidate_inf_value
% lean: CausalSmith.Stat.LdpAteEfficiencySurface.lowMean_Jstar_value

\item \textit{Uniqueness of the optimizing weights.}
Suppose that \(\beta\in\mathcal A_{\log3}\) satisfies
\[
\inf_{t\in\mathbb R}F_{\theta_{\mathrm{low}}}(\beta,t)
=
J^*(\theta_{\mathrm{low}},3/10,\log3).
\]
Its objective-value set is bounded below by zero, so the infimum is bounded above by its value at \(t_{\mathrm{low}}\). The dual upper bound then gives
\[
\frac{441}{3994}
=
\inf_{t\in\mathbb R}F_{\theta_{\mathrm{low}}}(\beta,t)
\le
F_{\theta_{\mathrm{low}}}(\beta,t_{\mathrm{low}})
\le
\frac{441}{3994}.
\]
All these values are equal. In particular, normalization of the dual sum yields
\[
\sum_{S\in\mathcal S}\beta_S f_S(\xi_{\mathrm{low}},t_{\mathrm{low}})
=
\frac{441}{3994}
=
\sum_{S\in\mathcal S}\beta_S
b_S(\log3)^\top\eta_{\mathrm{low}}.
\]
For an active pattern the corresponding terms are equal. For an inactive pattern, nonnegative multiplication of its strict slack gives
\[
\beta_S f_S(\xi_{\mathrm{low}},t_{\mathrm{low}})
\le
\beta_S b_S(\log3)^\top\eta_{\mathrm{low}}.
\]
Since these finitely many termwise inequalities have equal sums, equality holds in every term:
\[
\beta_S f_S(\xi_{\mathrm{low}},t_{\mathrm{low}})
=
\beta_S b_S(\log3)^\top\eta_{\mathrm{low}}
\quad(S\in\mathcal S).
\]
If \(S\notin\mathcal K\) and \(\beta_S\ne0\), feasibility implies \(\beta_S>0\). Its strict slack would then give
\[
\beta_S f_S(\xi_{\mathrm{low}},t_{\mathrm{low}})
<
\beta_S b_S(\log3)^\top\eta_{\mathrm{low}},
\]
a contradiction. Hence
\[
\beta_S=0\quad(S\notin\mathcal K).
\]

The normalization equations at input coordinates \(0,1,3\) now reduce to
\[
\begin{aligned}
3\beta_{S_1}+\beta_{S_6}+\beta_{S_8}&=1,\\
\beta_{S_1}+3\beta_{S_6}+\beta_{S_8}&=1,\\
\beta_{S_1}+\beta_{S_6}+3\beta_{S_8}&=1.
\end{aligned}
\]
Their differences give
\[
2(\beta_{S_1}-\beta_{S_6})=0,
\qquad
2(\beta_{S_6}-\beta_{S_8})=0,
\]
and substitution gives
\[
\beta_{S_1}=\beta_{S_6}=\beta_{S_8}=\frac15.
\]
For \(S\in\mathcal K\), these weights equal \(\alpha_S\); for \(S\notin\mathcal K\), both weights are zero. Therefore \(\beta=\alpha\).
% lean: CausalSmith.Stat.LdpAteEfficiencySurface.three_output_region

\item \textit{Reciprocal variance.}
Finally, the reciprocal definition in \cref{obj:def:staircase-saddle} and the positive optimized value give
\[
V^*(\theta_{\mathrm{low}},3/10,\log3)
=
\left(J^*(\theta_{\mathrm{low}},3/10,\log3)\right)^{-1}
=
\left(\frac{441}{3994}\right)^{-1}
=
\frac{3994}{441}.
\]
% lean: CausalSmith.Stat.LdpAteEfficiencySurface.three_output_region
\end{enumerate}
\end{proof}
